The one-particle space of the twistor attachment is built on $\CS$ (Theorem \claimref{field-theory.twistor-attachment}), and $\CS = \Lambda^{\mathrm{even}}W$ with the action of $\mathrm{Spin}(6)$, complexified $SL(4)$, by Theorem \claimref{rungs.two-twistor-bivector-null-invites-one}; $\Lambda^{\mathrm{odd}}W$ is the other half-spinor module, the dual $\CSm$. The element $-1 \in SL(4) = \mathrm{Spin}(6, \C)$ is the nontrivial element of the kernel of $\mathrm{Spin}(6) \to SO(6)$ and acts as $-1$ on both half-spinor modules, hence as $-1$ on $\CS$, on $\CSm$ and on all of $\Lambda^{\bullet}W$.

(1) The $n$-particle vectors $a^{\dagger}(u_{1})\cdots a^{\dagger}(u_{n})\Omega$ span the image of the equivariant map $(\C^{K}\otimes U)^{\otimes n} \to F$ (equivariant because $\Omega$ is invariant and $a^{\dagger}$ is equivariant and linear), so $F$ is a sum of quotients of the tensor powers. On $(\C^{K}\otimes U)^{\otimes n}$ the central element $z_{K} = (-1_{K}, 1)$ acts by $(-1)^{n}$, and so does $z_{4} = (1, -1_{4})$, since $-1_{4}$ acts by $-1$ on $U$ in each of the three cases. Hence $z_{K}z_{4}$ acts trivially on every tensor power and on $F$. If $K$ is even, $-1_{K} \in SL(K)$, so $z_{K}$ fixes every $SL(K)$-invariant vector, and then $z_{4}$ fixes it too. But $z_{4}$ acts by $-1$ on every vector of an $SL(4)$-isotypic component of type $\CS$ or $\CSm$. So such a vector is zero.

(2) In $\mathcal{S}$ the slot $SL(2)$ acts on $\Lambda^{j}\C^{2}$, trivially for $j = 0$ and $j = 2$ (the determinant is $1$), and $SL(4)$ acts on $\Lambda^{\bullet}W = \CS \oplus \CSm$; so $\mathcal{S}_{0}$ and $\mathcal{S}_{2}$ are $SL(2)$-invariant, each $\CS \oplus \CSm$, and $\nu^{c}, u^{c}$ sit at $(j, k) = (0, 0), (0, 2)$ and $e^{c}, d^{c}$ at $(2, 0), (2, 2)$ in Theorem \claimref{standard-model.under-s-composites-state-space-spinor}(1). An equivariant linear map carries $SL(2)$-invariant vectors of a given $SL(4)$-type to $SL(2)$-invariant vectors of the same type, which by (1) with $K = 2$ are zero in $F$; so it vanishes on $\mathcal{S}_{0} \oplus \mathcal{S}_{2}$. Dually, $z_{2}z_{4}$ acts trivially on $F$ and by $(-1)^{j}\cdot(-1) = -1$ on $\mathcal{S}_{0}\oplus\mathcal{S}_{2}$ and $+1$ on $\mathcal{S}_{1}$, so an equivariant map $F \to \mathcal{S}$ lands in the $+1$ eigenspace $\mathcal{S}_{1}$. An $SL(2)\times SL(4)$-equivariant subquotient of $F$ isomorphic to $\mathcal{S}$ would contain $z_{2}z_{4} = -1$ eigenvectors, impossible since $z_{2}z_{4} = 1$ on $F$ and on its subquotients.

(3) Computed by exact highest-weight counting. For $U = \CS$ realised as $\C^{4}$ with the standard $GL(4)$, an $SL(K)$-singlet copy of $\CS$ (of $\CSm$) is a $GL(K) \times GL(4)$ highest weight vector of weight $((m^{K}); (c+1, c, c, c))$ (respectively $(c+1, c+1, c+1, c)$) with $Km = 4c + 1$ (respectively $4c + 3$), and its multiplicity is the dimension of the joint kernel of the simple raising operators $E_{i,i+1} = \sum_{a}a^{\dagger}_{ia}a_{i+1,a}$ and $E_{a,a+1} = \sum_{i}a^{\dagger}_{ia}a_{i,a+1}$ on that weight space. For $K = 3$ the admissible weights in the exterior algebra (degree at most $12$) are $((1,1,1);(1,1,1,0))$ and $((3,3,3);(3,2,2,2))$, with six-dimensional weight spaces and multiplicity zero; in the symmetric algebra the first has multiplicity one (the vector $\sum_{\sigma}\mathrm{sgn}(\sigma)\,x_{1\sigma(1)}x_{2\sigma(2)}x_{3\sigma(3)}$ over bijections onto the first three twistor indices, the dual twistor of three constituents) and the second, with a $132$-dimensional weight space, multiplicity zero. For $K = 4$ no weight satisfies $4m = 4c + 1$ or $4c + 3$. For the exterior algebra and every $K \ge 2$ the same conclusion follows from skew Howe duality, $\Lambda(\C^{K}\otimes\C^{4}) = \bigoplus_{\lambda}S_{\lambda}\C^{K}\otimes S_{\lambda'}\C^{4}$ (imported, cited from memory): an $SL(K)$-singlet has $\lambda = (m^{K})$, and $S_{(K^{m})}\C^{4}$ is $K$ times the $m$-th fundamental weight, trivial for $m = 4$ and never $\CS$ or $\CSm$ for $K \ge 2$.

Remark. The obstruction needs $SL(4)$: no element of $U(2) \times U(3)$, acting as $\Lambda^{\bullet}(A \oplus B)$, acts as $-1$ on all of $\Lambda^{\bullet}W$, since $\Lambda^{0}W$ is fixed; for instance with $K = 2$ fermions and $\CS$ restricted to $U(3)$ as $\Lambda^{0}W \oplus \Lambda^{2}W$, $\Lambda^{2}(\C^{2}\otimes\CS) \supset \Lambda^{2}\C^{2}\otimes\mathrm{Sym}^{2}\CS \supset \det\otimes(\Lambda^{0}W\cdot\Lambda^{2}W) \cong \det\otimes\Lambda^{2}W$, the $d^{c}$-type component of $\mathcal{S}_{2}$ as a $U(2)\times U(3)$ module.

The script computes the highest-weight multiplicities of (3) and controls (one constituent: $\CS$ once in degree one; $K = 2$: the doublet copies of $\CS$ in degree one and of $\CSm$ at weight $((4,3);(2,2,2,1))$ found once each), checks that for $K = 2$ and $K = 4$ no target weight exists, and counts the sectors $\mathcal{S}_{0}$, $\mathcal{S}_{2}$ of $\Lambda^{\bullet}(\C^{2}\oplus\C^{3})$ by $W$-parity (four even and four odd states each).
